% !TEX root = ../MasterThesis.tex
\chapter{Streamlining CEGIS}
\label{chap:streamlining}

The preceding chapters characterised the two halves of certificate synthesis in
isolation: Chapter~\ref{chap:verification} the sound verifiers, and
Chapter~\ref{chap:refutation} the refuters as cheap counterexample search. This
chapter assembles them into the full CEGIS loop and asks: does placing a refuter in front of the sound verifier make
end-to-end certificate synthesis faster?

\section{Experimental Setup}

We run the complete CEGIS pipeline on three benchmark systems chosen to span
the difficulty spectrum of Chapter~\ref{chap:benchmarks}: the 2D $\mathrm{LinSwitch}$ (Section~\ref{sec:bench-linear}), the $\mathrm{DoubleWell}$ gradient
flow (Section~\ref{sec:bench-doublewell}), and
the LQR-controlled $\mathrm{Pendulum}$ (Section~\ref{sec:bench-pendulum-lqr}),
whose transcendental dynamics make bound propagation expensive. The fourth
benchmark, $\mathrm{LinStoch}$, is absent by necessity rather than choice:
Chapter~\ref{chap:verification} found it certifiable only by the exact MILP
route, and the LiRPA engines this campaign compares were unable to verify it,
so no verifier-only baseline exists to improve upon.

The learner trains the $\relu$ certificate architecture used throughout (two
hidden layers of eight units) to a training margin of $0.02$, and the loop
terminates when a sound verifier certifies the drift condition to
$\varepsilon = 10^{-3}$. Any counterexample found, whether by refuter or
verifier, is added to the training set and the candidate is retrained.

For our verifiers we take the three LiRPA-family engines of
Chapter~\ref{chap:verification}, the only family that certifies all three
systems. All three run inside the same branch-and-bound refinement with longest-edge region splitting. Against each engine we compare three loop configurations, which differ only in
how counterexamples are sought:

\begin{description}
\item[Verifier-only.] The baseline: the sound verifier runs every round, both
  finding counterexamples and, eventually, certifying.
\item[Random pre-screen.] Each round, random search
  (Section~\ref{sec:bbob}) hunts for a counterexample first; only when it
  finds nothing do we fall back on the sound verifier. Random is the
  uninformed extreme: essentially free, but blind.
\item[Whale pre-screen.] The same arrangement with the whale optimiser (WOA)
  at its BBOB-tuned hyperparameters (Table~\ref{tab:bbob-best}), selected because it dominated the fast-and-accurate
  frontier of the refuter benchmark.
\end{description}

Refuters run with the same budget as in Chapter~\ref{chap:refutation} ($8192$
evaluations per search, $16$ Monte-Carlo samples per point). Each
configuration is repeated over five seeds. We record the wall-clock time to a
verified certificate, the number of CEGIS rounds, the number of sound-verifier
calls, and the number of \emph{escapes} --- counterexamples the refuter
missed that the sound verifier subsequently found.

\paragraph{Compute environment.} Campaigns throughout the thesis were
dispatched as job arrays to a Slurm compute cluster; the sampling verifiers
(MC, MAB) and the three LiRPA engines benefit from GPU acceleration and ran on
GPU nodes. The MILP experiments are the exception, running locally on a
desktop CPU, as the Gurobi licence (version 13.0.1, via \texttt{gurobipy})
could not be reached from the cluster's compute nodes. All code (the
unified verifier framework, the refuters, the experiment drivers, and every
figure script) is available at
\url{https://github.com/will-bailkoski/streamlining-neural-certificates}.

\section{Results and Discussion}

%Table~\ref{tab:cegis-speedup} summarises the campaign.

\begin{table}[h]
\centering
\begin{tabular}{lllrrrrr}
\hline
Engine & System & Strategy & $n$ & Time (s) & Rounds & Calls & Speed-up \\
\hline
IBP & $\mathrm{LinSwitch}$     & verifier-only & 5 & $155.5$ & $43.8$ & $43.8$ & --- \\
    &                 & random        & 5 & $58.9$  & $18.0$ & $2.6$  & $2.64\times$ \\
    &                 & whale         & 5 & $39.8$  & $10.8$ & $1.0$  & $3.91\times$ \\
    & $\mathrm{DoubleWell}$     & verifier-only & 3 & $14.6$  & $1.3$  & $1.3$  & --- \\
    &                 & random        & 3 & $14.1$  & $1.7$  & $1.0$  & $1.03\times$ \\
    &                 & whale         & 5 & $18.4$  & $2.2$  & $1.0$  & $1.04\times$ \\
    & $\mathrm{Pendulum}$  & verifier-only & 1 & $55.0$  & $6.0$  & $6.0$  & --- \\
    &                 & random        & 2 & $33.7$  & $7.5$  & $1.0$  & $1.76\times$ \\
    &                 & whale         & 1 & $57.4$  & $10.0$ & $1.0$  & --- \\
\hline
CROWN & $\mathrm{LinSwitch}$     & verifier-only & 5 & $44.6$  & $11.8$ & $11.8$ & --- \\
      &                 & random        & 5 & $52.2$  & $15.4$ & $2.0$  & $0.85\times$ \\
      &                 & whale         & 5 & $38.5$  & $10.8$ & $1.0$  & $1.16\times$ \\
      & $\mathrm{DoubleWell}$     & verifier-only & 3 & $41.7$  & $1.7$  & $1.7$  & --- \\
      &                 & random        & 5 & $31.3$  & $2.2$  & $1.0$  & $1.48\times$ \\
      &                 & whale         & 5 & $32.0$  & $2.2$  & $1.0$  & $1.51\times$ \\
      & $\mathrm{Pendulum}$  & verifier-only & 3 & $203.1$ & $7.0$  & $7.0$  & --- \\
      &                 & random        & 4 & $57.8$  & $6.8$  & $1.0$  & $2.74\times$ \\
      &                 & whale         & 4 & $66.8$  & $8.5$  & $1.0$  & $3.19\times$ \\
\hline
$\alpha$-CROWN & $\mathrm{LinSwitch}$  & verifier-only & 5 & $935.0$  & $69.2$ & $69.2$ & --- \\
               &              & random        & 5 & $91.8$   & $15.4$ & $2.0$  & $10.19\times$ \\
               &              & whale         & 5 & $59.6$   & $10.8$ & $1.0$  & $15.70\times$ \\
               & $\mathrm{DoubleWell}$  & \multicolumn{6}{c}{\emph{(did not complete)}} \\
               & $\mathrm{Pendulum}$ & verifier-only & 4 & $9559.4$ & $8.0$  & $8.0$  & --- \\
               &              & random        & 4 & $1363.0$ & $6.8$  & $1.0$  & $5.97\times$ \\
               &              & whale         & 4 & $1331.5$ & $8.8$  & $1.0$  & $6.81\times$ \\
\hline
\end{tabular}
\caption{End-to-end CEGIS performance per engine, system and counterexample
strategy: mean wall-clock time, CEGIS rounds and sound-verifier calls over the
$n$ completed seeds. Speed-ups are ratios of seed-matched means against the
verifier-only baseline. $39$ of the $135$ configured arms did not complete
(hence the uneven $n$); in particular the $\alpha$-CROWN $\mathrm{DoubleWell}$ arms are
absent, and the IBP $\mathrm{Pendulum}$ cells are single-seed and admit no matched
comparison. The IBP $\mathrm{LinSwitch}$ baseline mean is inflated by a single runaway seed
(see text).}
\label{tab:cegis-speedup}
\end{table}

\textbf{Wall time improvements.} Stripping out the repeated sound calls changes
what a run spends its time on, not merely how long it takes. Averaged over the
three engines, sound verification is the largest single cost of a verifier-only
run --- a mean $62\%$ of its wall time, against $38\%$ for training and loop
overhead. A refuter pre-screen inverts this balance: verification falls to a mean
$47\%$ of wall time while training-and-overhead rises to $48\%$ and becomes the
dominant term, and the refuter's own search accounts for only about $5\%$ --- a
second or two per run at this budget. Once the refuter absorbs the counterexample
search, in other words, CEGIS is no longer verification-bound but
\emph{training}-bound.

How cleanly this holds depends on the engine's per-call cost. Under CROWN it is a
clean inversion --- verification drops from $57\%$ of a baseline run to $46\%$ and
retraining overtakes it at $49\%$. Under IBP the loop is training-bound from the
start: IBP's calls are so cheap that training is already $58\%$ of the
verifier-only run and merely edges up to $62\%$ with a refuter, so here the
refuter's payoff is rounds saved rather than call-time skipped. At the far end,
$\alpha$-CROWN's single sound call is costly enough that verification stays
dominant even after the repeats are gone --- still $68\%$ of a refuter run on
average, and $98\%$ on $\mathrm{Pendulum}$. Verification ceases to dominate exactly when a single sound
call is cheap relative to training, which is the regime the pre-screen is built to
create.

\textbf{Difficulty scaling.} On $\mathrm{Pendulum}$ under CROWN, where a single sound verification
costs $22$--$34$\,s, the refuter arms synthesise a certificate $2.7$--$3.2$
times faster than the baseline; on $\mathrm{LinSwitch}$ under $\alpha$-CROWN,
whose slope optimisation raises the per-call cost to $9$--$18$\,s, the speed-up
reaches $10.2\times$ for random and $15.7\times$ for whale --- the largest in
the campaign, because the verifier-only baseline there spirals to $69$ rounds
($935$\,s) that the refuter collapses to about $11$. Conversely, on the same
$\mathrm{LinSwitch}$ under plain CROWN a call costs well under a second, so there is
nothing worth saving: random search actually \emph{slows the loop down}
($0.85\times$), and whale's modest $1.16\times$ gain comes from needing fewer
rounds rather than from cheaper rounds. $\mathrm{DoubleWell}$ sits in between: most
seeds certify in one or two rounds, leaving little verification to skip. The
extreme case is $\mathrm{Pendulum}$ under $\alpha$-CROWN, where a single sound call
costs some $20$ \emph{minutes}: the verifier-only baseline makes six to eleven
of them --- one-and-a-half to three-and-a-half \emph{hours} per run --- while
each refuter arm makes exactly one and completes in around $22$ minutes, a
$6.8\times$ speed-up. At this price the pre-screen is the difference between an
overnight campaign and a coffee break.

The mechanism is visible in the calls column of Table~\ref{tab:cegis-speedup}.
The baseline pays one sound verification per round (up to seventy on the
$\alpha$-CROWN $\mathrm{LinSwitch}$), whereas the refuter arms finish with a handful
of sound calls at most, and whale with exactly one: the final, unavoidable
certification call. The refuters' own cost is negligible at this budget, a
second or two per run in total. The saving is therefore the per-call
verification cost multiplied by the rounds avoided, which is why the speed-up
tracks the engine and system so closely: skipping a dozen sub-second calls buys
nothing, while skipping several twenty-minute calls buys hours.

\textbf{Counterexample leakage.} In $12$ of $33$
completed random runs the refuter declared a candidate clean that the sound
verifier then refuted --- $18$ escaped counterexamples in all, each costing an
extra expensive verification round. Whale escaped nothing: in all $34$ of its
completed runs the first certificate it passed to the verifier was accepted,
so every whale run ended with exactly one sound call. Counterexample
\emph{quality} matters too, not just detection: on $\mathrm{LinSwitch}$ random
needed $15.4$ rounds on average against the baseline's $11.8$ while whale needed $10.8$,
matching the verifier-quality counterexamples of the baseline.

A reassuring consistency check: whale's round counts are identical, seed for
seed, under all three engines ($\mathrm{LinSwitch}$: $11,8,9,10,16$) --- because it
never lets the sound verifier find a counterexample, its loop trajectory
decouples from the verification engine entirely, whose only remaining role is
to price the final call. Random's trajectories match across CROWN and
$\alpha$-CROWN ($15.4$ rounds) but drift under IBP ($18.0$): each escape hands
control back to the engine, and a loose engine returns different, weaker
counterexamples.

\textbf{Engine choice.} Per-call costs span more than three orders of magnitude across the
campaign: IBP from $0.2$\,s ($\mathrm{LinSwitch}$) to $11$\,s ($\mathrm{DoubleWell}$), CROWN from
$0.5$\,s to $34$\,s, and $\alpha$-CROWN from $9$\,s to some $1200$\,s on
$\mathrm{Pendulum}$. Reaching for the cheapest engine instead of a refuter is a false
economy, however. IBP's loose enclosures produce weak counterexamples, and on
one $\mathrm{LinSwitch}$ seed the verifier-only loop spiralled: $179$ rounds, $633$\,s,
and a final certificate with a Lipschitz constant of $58$ against the $3$--$7$
typical of the other engines. Notably, the pre-screen \emph{rescues} this
failure --- fed on refuter counterexamples instead of IBP's own, the same seed
certifies in $16$ rounds ($56$\,s, $11\times$ faster) under whale. Here the
gain is not skipped call-time (IBP's calls are nearly free) but avoided rounds:
the refuter pays off through two distinct channels, skipping expensive sound
calls \emph{and} supplying better counterexamples than a loose verifier can
produce. A tight engine guarded by a refuter pre-screen obtains tight bounds
while paying the expensive call only once.

\section{Conclusions}

Sound verification is the price of a guarantee, but it need not be paid every
round. A cheap refuter in front of the verifier reduces CEGIS to a sequence
of fast search-and-retrain rounds followed by a single certification call,
and the resulting speed-up scales with how expensive that call is: from
nothing when verification is sub-second, to $3$--$4\times$ when calls cost
tens of seconds, to $15.7\times$ on the $\alpha$-CROWN $\mathrm{LinSwitch}$ runaway and a
practical necessity on the $\alpha$-CROWN $\mathrm{Pendulum}$, where $20$-minute sound
calls turn a multi-hour verifier-only run into a $22$-minute one. The benefit
flows through two channels: skipping
expensive sound calls, and feeding the learner better counterexamples than a
loose verifier can produce --- the latter rescuing IBP from a $179$-round
runaway. The quality of the refuter governs how much of this is realised: an
uninformed search leaks counterexamples back to the verifier and pads the
loop with uninformative rounds, while the tuned whale optimiser matched the
verifier's own counterexample quality at a fraction of its cost, never
leaked, and made the loop's trajectory independent of the engine beneath it.
The practical recipe for streamlining CEGIS is therefore to pair the tightest
verifier one can afford with a well-tuned refuter pre-screen, rather than to
economise on the verifier itself.

\section{Future Work}

The field offers no shortage of hard problems; we highlight the directions that
are opened directly by a finding of this thesis.

\paragraph{Nonlinear drift-Lipschitz bounds.} Both Lipschitz contributions of
Chapter~\ref{sec:lipschitz} lean on the linearity that keeps the drift
piecewise linear, but the principle behind them --- bound the object the
verifier consumes, rather than composing bounds on its factors --- does not.
Running bound propagation over the composed graph $\cert \circ f - \cert$
would extend the tight drift constant, and its per-axis refinement, to
nonlinear dynamics.

\paragraph{Closing the $\mathrm{LinStoch}$ gap.} The capability boundary of
Chapter~\ref{chap:verification} leaves exactly one system uncertified by any
relaxation-based engine: $\mathrm{LinStoch}$, whose noise-floor equilibrium
leaves the relaxations too little margin to work with. Diagnosing that
interaction (for instance by decoupling the equilibrium the learner trains
against from the one the verifier certifies) is a self-contained problem
whose solution would make the bound-propagation family genuinely
suite-spanning.

\paragraph{Rescuing sound sampling.} The sampling verifiers of
Chapter~\ref{chap:verification} failed on loose inputs rather than broken
theory: termination stalled on the magnitude of the Lipschitz and statistical
radii. The anisotropic per-axis constants of Chapter~\ref{sec:lipschitz} plug
directly into the bandit's splitting rule, and variance-adaptive concentration
bounds attack the statistical term. Success here would be significant, as it
would deliver a sound verifier for genuinely black-box dynamics, which no
other family can offer.

\paragraph{Refuter--verifier co-design.} The pre-screen of this chapter treats
the refuter and the verifier as strangers who share only a training set. A
natural next step is to let the refuter's near-violations warm-start the
verifier's branch-and-bound, seeding splits where the margin is thinnest,
and to make the fall-back policy adaptive, spending refuter budget in
proportion to how often it has recently let a counterexample escape.

\paragraph{Broader specifications and certificate classes.} We claimed in
Chapter~\ref{sec:preliminaries} that our methods extend to safety and
reach-avoid specifications without essential modification; that claim deserves
an experiment. Likewise, this thesis used piecewise-linear certificates
throughout: smooth activations such as $\tanh$ remove the symbolic engines
from play but remain within reach of bound propagation, and
Lipschitz-constrained architectures interact directly with the margin budget
\eqref{eq:budget}. Both would re-draw the capability boundary this thesis
mapped.


\newpage

